% 米氏散射：瓊斯矩陣、斯托克斯向量與穆勒矩陣（骨架；完整推導見 Bohren & Huffman (1983) 第二至四章）

\subsubsection*{(六) 遠場與振幅散射矩陣（瓊斯矩陣）}
本節只列出結果，推導見 \textcite{bohren1983Absorption} 第三、四章；瓊斯矩陣的寫法源自 \textcite{jones1941NewCalculus}。

在遠場（$kr \gg x^2$），散射場為橫向的局部平面波。以包含入射方向與散射方向的平面為散射面，將電場分解為平行（$\parallel$）與垂直（$\perp$）於散射面的分量，散射場與入射場的關係為
\begin{equation}
    \begin{pmatrix} E_{s\parallel} \\ E_{s\perp} \end{pmatrix} = \frac{e^{ikr}}{-ikr} \begin{pmatrix} S_2(\theta) & 0 \\ 0 & S_1(\theta) \end{pmatrix} \begin{pmatrix} E_{\parallel i} \\ E_{\perp i} \end{pmatrix} , \label{eq:jones}
\end{equation}
其中振幅散射函數為
\begin{equation}
    S_1(\theta) = \sum_{n=1}^{\infty} \frac{2n+1}{n(n+1)} \left( a_n \pi_n + b_n \tau_n \right), \qquad S_2(\theta) = \sum_{n=1}^{\infty} \frac{2n+1}{n(n+1)} \left( a_n \tau_n + b_n \pi_n \right) . \label{eq:S1S2}
\end{equation}
球對任何通過球心的軸旋轉對稱，因此矩陣的非對角元為零，且 $S_1$、$S_2$ 只與散射角 $\theta$ 有關。$\parallel$ 與 $\perp$ 分量各自乘上 $S_2$、$S_1$；一般 $S_1 \neq S_2$，所以散射會改變偏振態。

\subsubsection*{(七) 斯托克斯向量}
斯托克斯參數由 \textcite{stokes1852Composition} 提出，本節的定義與參考面旋轉依 \textcite{bohren1983Absorption} 第二章。斯托克斯參數定義為
\begin{equation}
    \begin{aligned}
        I &= \langle E_{\parallel} E_{\parallel}^* + E_{\perp} E_{\perp}^* \rangle, \qquad & Q &= \langle E_{\parallel} E_{\parallel}^* - E_{\perp} E_{\perp}^* \rangle \\
        U &= \langle E_{\parallel} E_{\perp}^* + E_{\perp} E_{\parallel}^* \rangle, \qquad & V &= i \langle E_{\parallel} E_{\perp}^* - E_{\perp} E_{\parallel}^* \rangle
    \end{aligned} \label{eq:stokes_def}
\end{equation}
其中 $\langle\cdot\rangle$ 為時間平均，斯托克斯向量為 $\mathbf{S} = (I, Q, U, V)^{T}$；完全偏振光滿足 $I^2 = Q^2 + U^2 + V^2$。若把參考軸繞傳播方向轉動角度 $\alpha$，$I$、$V$ 不變，$Q$、$U$ 依下式轉換：
\begin{equation}
    \mathbf{S}' = \mathbf{R}(\alpha)\, \mathbf{S}, \qquad \mathbf{R}(\alpha) = \begin{pmatrix} 1 & 0 & 0 & 0 \\ 0 & \cos 2\alpha & \sin 2\alpha & 0 \\ 0 & -\sin 2\alpha & \cos 2\alpha & 0 \\ 0 & 0 & 0 & 1 \end{pmatrix} . \label{eq:rot}
\end{equation}

\subsubsection*{(八) 穆勒矩陣}
本節推導依 \textcite{bohren1983Absorption} 第三、四章；在蒙地卡羅模擬中逐次「先旋轉參考面、再乘穆勒矩陣」的做法見 \textcite{ramellaRoman2005ThreeMC} 與 \textcite{bartel2000MonteCarlo}。

將式\eqref{eq:jones}代入式\eqref{eq:stokes_def}，得散射面座標下散射光與入射光斯托克斯向量的關係：
\begin{equation}
    \mathbf{S}_s = \frac{1}{k^2 r^2} \mathbf{F}(\theta)\, \mathbf{S}_i, \qquad \mathbf{F} = \begin{pmatrix} F_{11} & F_{12} & 0 & 0 \\ F_{12} & F_{11} & 0 & 0 \\ 0 & 0 & F_{33} & F_{34} \\ 0 & 0 & -F_{34} & F_{33} \end{pmatrix} , \label{eq:mueller}
\end{equation}
\begin{equation}
    \begin{aligned}
        F_{11} &= \tfrac{1}{2}\left( |S_1|^2 + |S_2|^2 \right), \qquad & F_{12} &= \tfrac{1}{2}\left( |S_2|^2 - |S_1|^2 \right), \\
        F_{33} &= \mathrm{Re}\left( S_2 S_1^* \right), \qquad & F_{34} &= \mathrm{Im}\left( S_2 S_1^* \right) .
    \end{aligned}
\end{equation}
$\mathbf{F}$ 即球的穆勒矩陣，只有四個獨立元素。非偏振光入射時，散射強度正比於 $F_{11}(\theta)$，因此 $F_{11}(\theta)/\int F_{11}\,d\Omega$ 可作為蒙地卡羅模擬抽樣散射角的相位函數。

穆勒矩陣只在散射面座標下具有式\eqref{eq:mueller}的形式。若入射光的斯托克斯向量 $\mathbf{S}_{lab}$ 以實驗室參考軸表示，須先旋轉到散射面再乘 $\mathbf{F}$：
\begin{equation}
    \mathbf{S}_s = \frac{1}{k^2 r^2}\, \mathbf{F}(\theta)\, \mathbf{R}(\alpha)\, \mathbf{S}_{lab} ,
\end{equation}
其中 $\alpha$ 為實驗室參考軸轉到散射面所需的角度。蒙地卡羅模擬每次散射都依此「先旋轉、再乘 $\mathbf{F}$」更新光子的斯托克斯向量。
